\section{Il Capo}

\begin{quotex}
The rule of many is not good; one ruler let there be.
\flright{Ulysses, in Homer's \textit{The Iliad}, and again by Aristotle at the end of Book XII of the \textit{Metaphysics.}}
\end{quotex}


Much is made over the alleged Evolian reversal of the roles of the castes. This in part stems from his portrayal of the regal initiation being above the sacerdotal initiation, which arises from a misreading of some Hindu texts as Coomaraswamy pointed out. Evola also objected to the necessity for the Holy Roman Emperor to be consecrated by the Pope. However, the consecration is not the problem, but the idea of two rulers is. As Guenon points out the \emph{The Lord of the World}, the idea of a single head is certainly traditional and the split in the Middle Ages was a deviation. Guenon writes:

\begin{quotex}
By the Middle Ages, supreme power had already become divided between the Empire and the Papacy. Such a division marks an organisation that is incomplete at its head since the common principle, on which the two powers depend, is missing.
\end{quotex}


So with the intent of completely putting this issue to rest, we will bring out the full meaning of Guenon's point, and with a minor tweak, this will bring Evola completely in line with Tradition on this point. In \emph{La Tradizione Romana} Guido de Giorgio dedicated a section of the book to the ``Constitution of a Traditional Society''. The first three chapters deal with the three twice-born castes and the fourth with the idea of the ``Capo'' --- Head, Chief, Leader. Beginning with Ulysses' statement, de Giorgio makes these points.

\begin{itemize}


\item A single Chief in time corresponds to the supreme unity in eternity.

\item Since God is pure contemplation, so inversely the Chief will make of his life a pure activity dedicated to the preservation of His order in the world.

\item The Chief holds temporal power and supreme authority whose dominion embraces all of active life.

\item His work depends on the highest virtue: Justice.

\item He the the Prince of the Warriors, but his power is exercised over everyone.

\item Although he \emph{arises} from the Warrior caste, he is \emph{above} all the castes.

\item He is consecrated by the first caste.

\item He leads the second caste in the protection and preservation of Traditional order.

\item He guards the ``Temple'' and protects the contemplatives so that spiritual truths are available to all.

\item He is absolutely autonomous and the priests own him obedience, just as do the other castes.

\item The priests must not interfere in temporal affairs.
\end{itemize}


To summarize: although the Chief is consecrated by the priests, his power is by Divine Right. He holds temporal power over all. His is a life of activity, since he is the means through which God's will ``Providence'' is accomplished. In other words, by his dominion and action, he does not experience privation.


\flrightit{Posted on 2011-03-11 by Cologero}

\begin{center}* * *\end{center}

\begin{footnotesize}
\begin{sffamily}
\texttt{Albion on 2011-03-11 at 01:09 said:}

Could one say that there is a difference between ecclesiastical/bureaucratic unity and visible unity?

\hfill

\texttt{Visionsofglory14 on 2011-03-11 at 02:09 said:}

Evola seems to question the necessity of consecration. To him, it only becomes necessary in later, more decadent periods when sacerdotal temporal power begins to encroach on the regal. The true Brahman, in Evola's interpretation, would not need consecration since his spiritual authority speaks for itself. Divine Right cannot be granted except by the spirit itself, so in his highest formation, the Chief would not need consecration.

\hfill

\texttt{Albion on 2011-03-11 at 08:00 said:}

So what we would normally conceive as emerging in and from the invisible world to ``manifest'', Evola considers to have had an actual physical emergence in its own invisibility? For those who could see?

\hfill

\texttt{Ted on 2011-03-11 at 18:03 said:}

But for what it is worth: Evola does cite the Brihad Aranyaka Upanishad, in which a Kshatriya instructs a Brahmin, and when the Brahmin expresses surprise at his knowledge he replies that there were some secrets that the Kshatriyas still have not told the Brahmins. This is connected with Tantra, by the way. I don't have time to go into this right now. Perhaps tommorow.

One last thing: In Volume One of Will Durant's ``The Story of Civiliztion''(``Our Oriental Heritage''), he asserts that orginally the Kshatriya class was superior to the Brahmins, who were their assistants. Ananda Coommarasawmy was one of the Durant's main consultants on this volumne.

\hfill

\texttt{Cologero on 2011-03-11 at 19:24 said:}

We have always written about the castes in terms of their roles and the type of persons of which they consist. This has nothing to do with ``superiority'' or ``inferiority'', except only incidentally. The role of the Kshatriya is to rule and administer. If that is what you mean by ``superior'', it really adds nothing.

\hfill

\texttt{kadambari on 2011-03-11 at 20:51 said:}

... While Kings could have contempt for Brahmins as Kings, they also feared their knowledge and purity, and many sages had contempt for the life of the ruler or administrator and were feared and admired due to their wisdom, and even Kings feared the sages... In the Politics, Aristotle places the life of contemplation above the political life which is interesting, the enlightened man goes into it by a kind of necessity...

\texttt{HOO on 2011-03-11 at 21:21 said:}

Aristotle is writing very late in the transience of this planet (writing itself being a late degeneration), the split already being quite antique in his age; his perspective thus reflecting that.

\hfill

\texttt{kadambari on 2011-03-11 at 21:51 said:}

The Greek term ``Deinos'' captures the respect and awe which superior beings evoked... 

\hfill

\texttt{Ted on 2011-03-12 at 11:55 said:}

I hesitate to go into this... but here goes anyway.

The Hyperboreans. A divine or semi-divine race which lived within the northern circumpolar region. Due to glaciation brought on by physical or meta-physical reasons, or perhaps both, they gradually moved south. In India, they were known as the Aryans.

In Genesis Chapter 6, as well as in the Book of Enoch, there is mention of the ``Nephilim'' or the ``Watchers''. Without going into detail here, which in any event can by found Evola's ``Revolt Against the Modern World'', this myth refers to the interbreeding of this race with the local populations with which they came into contact with. The progeny were the Nephilim or the Watchers, and they had superior abilities, both physical and metaphysical, than did the common run of ``humanity''. As Plato says in the Critias, when Posiedon mated with human woman on Atlantis, the progeney which resulted remained pure and noble until the human admixture became predominant. At that point they became corrupt. And Atlantis was destroyed, either through war, technological catastrophe, the misuse of magic or by an act of higher powers. The underlying cause of this destruction, of course, was the corruption of its population.

In India, a genetic analysis of the population conducted 7 or 8 years ago showed that the Kshatriya class in India had the highest proportion of foreign DNA of a European origin. The Brahmins were a somewhat close second. And the other two castes were far distant.

This seems consistent with Evola's contention that the Kshatriya class was the superior; the Hyperborean/Aryans set up those most close to the original genetic stream, the Kshatriyas, as the rulers, and then selected the close second Brahmins to be their assistants.

\hfill

\texttt{Cologero on 2011-03-12 at 13:14 said:}

Wow, Ted, life would be so much easier in your sand-castle in the sky. Instead of having to prove his valor, intelligence, strength, virtue or leadership skills, all a man would have to do is display his genome and the masses of his genetic inferiors would make him dictator by popular acclaim.

And it is remarkable that men with high proportion of European DNA were or are incapable of performing the tasks appropriate to ``lower'' castes such as medicine, architecture, law, scientific research, art, and so on.

You obviously use Evola to support your claim (show us where he said the Kshatriya is superior because of his genetic stream); you may want to supplement this comment with what Evola says about the race of the spirit and the race of the body.

And explain to us, Ted, how these little factoids lead up to a spiritual transformation. Or should we prefer to sit around in despair cursing ourselves for not having chosen better parents?

\hfill

\texttt{Ted on 2011-03-13 at 12:38 said:}

My purpose, of course, was not to describe the entire Indian caste system. I was merely discussing the rank of the various castes in light of esoteric teaching and genetic evidence.

\hfill

\texttt{Cologero on 2011-03-13 at 14:25 said:}

Unfortunately, Ted, Evola never said what you say he said, just the opposite, in fact.

In the Chapter ``Regression of Castes'' in Revolt, he writes, regarding the process of regression:

``Authority descended to the level immediately below, that is, to the caste of the warrior.''

Hence, the warrior caste was never the first caste.

As for the genetic component, Evola specifically rejects this:

``regality of blood replaced regality of the spirit.'' (which he considers a negative development)

Please end this discussion.

\hfill

\texttt{Ted on 2011-03-13 at 16:11 said:}

But Evola does not, of course, reject the Hyperborean thesis. The influence of genetics and of the blood is transmitted to the descendents of the Hyperboreans, after their coupling with mortal women. Is this not a key component of traditonal teaching?

\texttt{kadambari on 2011-03-17 at 13:49 said:}

@Graham

I am not one to believe that everything was perfect in traditional societies. For example, in Persia, the religion was an aristocratic religion, and many of the peasantry disliked the priests and thought them corrupt. But that does not mean that they could not sort out their affairs and needed to get all converted by the the sword to an alien faith. It goes without saying that Persia makes Islam palatable through Sufism as it could not get out of it, changing its nature, Persianizing it; the real Isalm can be found in Saudi and Pakistan, not the Persianized Islam... 

In India, unlike Persia, movements like Buddhism lead to spreading of literacy. Buddhism is an internal pehenomenon and sought to make tradition stronger by seeking to strengthen those aspects which were getting corrupted. You also have religions like Jainism in which the mercantile classes who perhaps felt were left out could express spirituality in a way proper to their culture, and so on. But none of these religons destroy the essetial Indic unity of culture, which is why all these religions exist pretty much harmoniously. A culture can self-correct if it is left to itself, but when damage is done by external forces to an extent that an entirely way of life is lost, then it is exceedingly difficult to recapture that which is lost.

In India, the destruction was so great after the invasions that people's lives were destroyed and they would just wander about aimlesslessly after the devastations as they have no clue as to how to react as something they cannot understand has been imposed, this incidentally is where thugs come from, from people wandering about aimlessly after the invasions who have been rendered destitute who trun into thugs... 

A proper undestanding of Indic religions requires study at a deep level as there are many layers and nuances. My take on it is that if anything good survives there, it is nothing short of miraculous, as other societies never recovered from such devastations... 

So the process of healing is likely to be a gradual one if it does not get subverted by forces that prevent healing of a culture, most of which today is home grown so Indians of course are fully responsible. I am just saying there is a larger backdrop to the state of affairs there, the history of the region is a complex one... 

\hfill

\texttt{Cologero on 2011-03-18 at 13:19 said:}

There were 3 editions of Revolt, originally published in 1934, then revised and enlarged in 1951, with a final revision in 1969. I have the edition by Edizioni Mediterranee, which is a reprint of the third edition.

HOO, I only mentioned Gregor in regard to Evola, not to his work in general. Despite Gergor's important works on the actual political philosophy of Fascism, he took pains to exclude Evola as a serious thinker in that vein. You may care to look up what we wrote about Fabre d'Olivet who did in fact write about the conflict between the Hyperboreans and black races to the South.

The issue at stake, HOO, which both you and Ted are both unable to grasp, is the false conflation of race and caste. Nothing written here, nor by Evola and Guenon, would support that. You yourself alluded to the work of Dumezil who recognized a tri-functional structure, not tri-racial. Thus in the Hittite civilization you mention. Hittites are members of all three castes, there is no racial segregation implied in that in the least. Do not confuse the caste system with an apartheid system; the latter is racially based, with two parallel and separated cultures occupying the same geographical space. This is not what the caste system is.

You quote Evola: ``It is not a matter of indifference that a body has this shape rather than that one: it is not by chance and without consequence.''

Exactly our point. It is not a matter of chance since it is a reflection of the spirit, not the cause of the spirit. Genetic evolution, by definition, is based on chance and random variations. How can that account for a spiritually superior race? Now Evola also wrote that if a man is not twice-born, then he is no different from any other man, whatever the color of his eyes. He is just at the first stage that is described in Individual and Becoming. Digest that for a day or two before I provide the full quote.

You can criticize the Dravidians, but they have at least kept the Vedic and Tantric Tradition alive, whereas the blue-eyes have not. There is no serious (i.e., not sensationalized by Hollywood) portrayals of the Illiad, Beowulf, the Eddas, and so on. This is fundamentally a spiritual problem, not a biological issue. This must be seen and understood, or no progress is possible.

\hfill

\texttt{John on 2011-03-29 at 19:33 said:}

I was brought up to believe that the ideal was the warrior-scholar. This ideal is rarely or ever attained. Usually what you get is either an intelligent, knowledgeable, some-what cultured warrior or a scholar with a strong, stead-fast character.

In some families, rather than scholars giving birth to scholars, or warriors giving birth to warriors, you have scholars giving birth to warriors and warriors giving birth to scholars.

One thing that makes Evola attractive is that he is a scholar with a warrior's temperament. I know nothing about his father, but perhaps his father was a warrior.

In families that produce both warriors and scholars, then one or the other temperament will rule in a given temperament, but the other temperament will be that of his father, and therefore the temperament to which he owes obedience.

\hfill

\texttt{Cologero on 2011-03-29 at 19:55 said:}

Good point, John, the elusive philosopher-king. Perhaps the Templars embodied that ideal, maybe Marcus Aurelius or Frederick II. But even if that idea is attained in one individual, it never alters the fundamental subservience of temporal power to spiritual authority.

\end{sffamily}
\end{footnotesize}

